\section{The original ample line on the actual Frobenius twist}

Let $f:X\to\operatorname{Spec}k$ be the original structural morphism
in prime characteristic $p$.  Write $X'=X\times_{k,F_k}k$, let
$\pi:X'\to X$ be the first projection, and let
$F_{X/k}:X\to X'$ be relative Frobenius.  These maps satisfy the
actual scheme equality $\pi\circ F_{X/k}=F_X$.

\begin{definition}[The twisted original line]
\label{def:kn-frobenius-twisted-line}
\leanok
\lean{EnvelopeProperty.Vanishing.FrobeniusLineComparison.twistedLine}
For an original module $L$ on $X$, its twisted module is
$L'=\pi^*L$ on the actual scheme $X'$.
\end{definition}

\begin{lemma}[The relative pullback is the original absolute pullback]
\label{lem:kn-relative-line-pullback}
\leanok
\lean{EnvelopeProperty.Vanishing.FrobeniusLineComparison.relativePullbackIso}
\uses{def:kn-frobenius-twisted-line}
There is a canonical isomorphism
$F_{X/k}^*L'\simeq F_X^*L$.  No perfectness, line-frame, or smoothness
assumption is required for this comparison.
\end{lemma}
\begin{proof}
The pullback-composition isomorphism identifies the left-hand module
with $(\pi\circ F_{X/k})^*L$.  Transport along the actual scheme equality
$\pi\circ F_{X/k}=F_X$.  This gives the claimed isomorphism on the original
modules rather than a comparison of chosen affine presentations.
\end{proof}

\begin{lemma}[The same comparison retains every tensor power]
\label{lem:kn-relative-line-power-pullback}
\leanok
\lean{EnvelopeProperty.Vanishing.FrobeniusLineComparison.relativePullbackPowerIso}
\uses{lem:kn-relative-line-pullback}
For every $m\geq0$, there is a canonical isomorphism
$F_{X/k}^*((L')^{\otimes m})\simeq(F_X^*L)^{\otimes m}$.
\end{lemma}
\begin{proof}
Pullback commutes with the recursively defined actual tensor powers.
Apply the tensor power of the preceding module isomorphism.  At
$m=0$, both comparisons are the structure-module pullback comparison.
\end{proof}

\begin{lemma}[Absolute Frobenius raises original transition scalars to the $p$th power]
\label{lem:kn-frobenius-unit-transition}
\leanok
\lean{EnvelopeProperty.Vanishing.FrobeniusLineComparison.pullback_unitScalar,EnvelopeProperty.Vanishing.FrobeniusLineComparison.pullback_unitAutomorphism}
For an original global scalar $a$ on $X$, conjugating the pullback of
its structure-module multiplication map by the canonical unit
comparison gives multiplication by $a^p$.  For a unit $a$, this is an
equality of actual module automorphisms.
\end{lemma}
\begin{proof}
Pullback of scalar multiplication is multiplication by the actual
pullback of that scalar.  The original scheme Frobenius acts on every
original section by its $p$th power.  Substitution yields the scalar
equality.  For units, the same equality for forward maps proves equality
of their module isomorphisms.
\end{proof}

The additional isomorphism $F_X^*L\simeq L^{\otimes p}$ for an
arbitrary original line requires the equality of the full pulled and
tensor-power frame transitions and their geometric gluing.  The scalar
identity above controls frame changes; matching transitions and descent
will give that global isomorphism.

\begin{lemma}[The local Frobenius/power comparison is independent of frame]
\label{lem:kn-frobenius-frame-independence}
\leanok
\lean{EnvelopeProperty.Vanishing.FrobeniusLineComparison.powerFrame_hom_eq_comp_scalar,EnvelopeProperty.Vanishing.FrobeniusLineComparison.frameComparison_eq}
\uses{lem:kn-frobenius-unit-transition}
On an open where $L$ has an actual rank-one frame, compare $F_X^*L$
and $L^{\otimes p}$ by their induced frames.  The resulting module
isomorphism is independent of the chosen original rank-one frame.
\end{lemma}
\begin{proof}
Any two rank-one frames differ by multiplication by a unique unit $a$.
The actual pulled frames differ by multiplication by $a^p$, by the
preceding lemma.  By induction on the tensor exponent, the actual
tensor-power frames differ by multiplication by the corresponding
power of $a$: tensoring changes by $r$ and $s$ changes the tensor frame
by $rs$, since the structure-module tensor comparison is multiplication.
In degree $p$ the transition is therefore also $a^p$.  These common
invertible transitions cancel in the coefficient-preserving comparison.
\end{proof}

\begin{lemma}[The transported Frobenius frame respects further pullback]
\label{lem:kn-frobenius-frame-restriction}
\leanok
\lean{EnvelopeProperty.Vanishing.FrobeniusLineComparison.frobeniusFrame_factor,EnvelopeProperty.Vanishing.FrobeniusLineComparison.localComparison_eq}
\uses{lem:kn-frobenius-frame-independence}
The Frobenius frame transported from a genuine frame over $g:Y\to X$
and then pulled further through $j:Z\to Y$ agrees with the Frobenius
frame transported from the original frame pulled along $j$.  The
coefficient-preserving local comparison is independent of that
original frame.
\end{lemma}
\begin{proof}
Pullback composition identifies both frame maps with the pullback of
the original frame along the same composite morphism to $X$.
The Frobenius naturality square for $j$ identifies $F_Y\circ j$
with $j\circ F_Z$, and the one for $g$ identifies $F_X\circ g$
with $g\circ F_Y$.  Their composite is the Frobenius naturality
square for $g\circ j$.  Compatibility of the native
composition and structure-module comparisons therefore identifies
the two entire frame maps.  Applying frame independence on the
source scheme proves the second assertion.
\end{proof}

\begin{theorem}[The original Frobenius pullback is the original line power]
\label{thm:kn-frobenius-line-power}
\leanok
\lean{EnvelopeProperty.Vanishing.FrobeniusLineComparison.absolutePowerIso,EnvelopeProperty.Vanishing.FrobeniusLineComparison.relativePowerIso}
\uses{lem:kn-frobenius-frame-restriction,lem:kn-relative-line-pullback,thm:scheme-original-frame-gluing}
For every original invertible module $L$,
$F_X^*L\simeq L^{\otimes p}$ and
$F_{X/k}^*L'\simeq L^{\otimes p}$ on the same original scheme.
\end{theorem}
\begin{proof}
Take the genuine trivializing open cover of $L$.  On each open the
actual Frobenius naturality square transports its line frame to a
frame of $F_X^*L$.  The tensor-power construction gives a frame of
$L^{\otimes p}$ on that same open.  Their coefficient-preserving
comparison is independent of the original frame by the preceding
lemma.  Restriction of the transported Frobenius frame agrees with
transport of the restricted original frame: both are the pullback of
the original frame along the same Frobenius commutative triangle,
using composition compatibility of native pullback.  The same
composition compatibility for tensor-power pullback gives this
restriction property for the tensor-power frame.  Consequently the
local comparisons agree on every overlap.  Sheaf descent of actual
module morphisms glues these comparisons and their inverses to a
global module isomorphism.  Compose with the relative-to-absolute
comparison to obtain the relative Frobenius assertion.
\end{proof}


\begin{corollary}[Projection retains the original amplified tensor term]
\label{cor:kn-frobenius-line-projection}
\leanok
\lean{EnvelopeProperty.Vanishing.FrobeniusLineComparison.relativeTensorPowerIso,EnvelopeProperty.Vanishing.FrobeniusLineComparison.moduleProjectionIso}
\uses{thm:kn-frobenius-line-power,thm:fc-original-scheme-line-projection-formula}
For every original module $M$ and every $m\geq0$,
\[
 F_{X/k*}M\otimes (L')^{\otimes m}
 \simeq F_{X/k*}(L^{\otimes pm}\otimes M).
\]
\end{corollary}
\begin{proof}
The projection formula for the invertible module $(L')^{\otimes m}$
gives the pushforward of $M\otimes F_{X/k}^*((L')^{\otimes m})$.
Pullback commutes with tensor powers, and the preceding theorem
identifies that pulled line with $(L^{\otimes p})^{\otimes m}$.
Tensor-power multiplication identifies it with $L^{\otimes pm}$.
The symmetry of the original module tensor product gives the stated
order of factors.  This is a module comparison in each degree; the
differential is transported from the twisted pushed complex and is
not obtained by pretending the original de Rham differential is
$\mathcal O_X$-linear.
\end{proof}
